\SourceLeaf{337}%
\(T^0\) may be used), apply the four formulas
\[
\begin{aligned}
\beta&=(2)r\cos T\\
B&=B^0-\tfrac12\beta\\
\tau&=(1)r\sin T\operatorname{tang}B\\
T&=T^0-\tfrac12\tau
\end{aligned}
\]
The value of (2) should be taken from the table with \(B\) as the argument. For the first calculation one can use \(B^0\), unless one thinks a closer value of \(B\) can be estimated in advance. Take the value of (1) from the table using the value of \(B\) just found.

Repeat the same calculation with the value of \(T\) obtained from the fourth equation, taking (1) and (2) from the table using the already improved \(B\). Repeat again if necessary until the result stabilizes, that is, until the fourth formula returns the same value of \(T\) that was last used as the starting point. Use only five-figure logarithms for all these calculations.

In subsequent repetitions, use seven-figure logarithms, include the logarithmic corrections to \(\log\tau\) and \(\log\beta\), and set \(B=B^0-\tfrac12b,\ T=T^0-\tfrac12t\). Only once this calculation has also given stable results should \(\lambda\) and \(l\) be calculated from the formulas at the end of Article 33. To illustrate these instructions, here are the main steps of an example based on the side Brocken-Inselsberg of the triangle, just as in the spherical calculation in Article 20 above.

In the calculation for the ellipsoid, the latitude of Brocken is
\[
=51^0\ 48'\ 1''\ 9294=B^0
\]
and the azimuth of the side Brocken-Inselsberg is
\[
=5^0\ 42'\ 21''\ 7699=T^0
\]
The logarithm of the side of the triangle in toises is the same to the seventh decimal place as in the conformal representation on the sphere, namely \(=4{,}7353929\). Hence, in the unit on which our auxiliary table is based, \(\log r=5{,}0251757\).
